% ============================================================
% Capítulo 8 — TextReIDNet: conceptos + código
% ============================================================
\chapter{TextReIDNet: joint embedding visual-textual}
\label{cap:textreidnet}

Este capítulo explica los conceptos detrás del modelo principal y
muestra su código línea por línea.

\section{¿Qué es un joint embedding?}

\subsection{El concepto}

Un \emph{joint embedding space} (espacio de embedding conjunto)
es un espacio vectorial donde las dos modalidades --- imágenes
y texto --- se mapean al mismo lugar. La idea clave:

\begin{itemize}
  \item Imagen de una persona $\to$ vector $\mathbf{I}' \in
        \mathbb{R}^{1024}$.
  \item Descripción de esa persona $\to$ vector $\mathbf{T}' \in
        \mathbb{R}^{1024}$.
  \item \textbf{Si el modelo es bueno, los vectores están
        \emph{cerca} en el espacio}.
\end{itemize}

\noindent\textbf{Entonces} medir qué tan "cerca" están dos
vectores nos dice qué tan parecidas son las dos entradas.

\section{¿Qué es la similitud coseno?}

Dados dos vectores $\mathbf{a}, \mathbf{b} \in \mathbb{R}^n$:

\begin{equation}
  \text{sim}(\mathbf{a}, \mathbf{b}) = \frac{\mathbf{a} \cdot \mathbf{b}}
                                            {\|\mathbf{a}\| \cdot \|\mathbf{b}\|}
  \in [-1, +1]
  \label{eq:cosine}
\end{equation}

\noindent\textbf{Interpretación}:
\begin{itemize}
  \item $+1$: vectores idénticos (apuntan en la misma dirección).
  \item $0$: ortogonales (no relacionados).
  \item $-1$: opuestos.
\end{itemize}

\noindent\textbf{Por qué coseno y no distancia euclidiana?}:
\begin{itemize}
  \item Invariante a la magnitud: $\mathbf{a}$ y $2\mathbf{a}$
        tienen similitud 1.
  \item Es un \emph{ángulo} entre vectores, no una distancia.
  \item Valores acotados en $[-1, +1]$, fácil de interpretar.
\end{itemize}

\section{¿Por qué un DSC compartido?}

La última capa (\texttt{DepthwiseSeparableConv(1024, 1024)}) se
usa para \textbf{ambas} ramas (visual y textual). Esto tiene dos
efectos:

\begin{enumerate}
  \item \textbf{Menos parámetros}: una capa en lugar de dos.
  \item \textbf{Mismo espacio}: las dos modalidades se proyectan
        con la misma transformación, lo que \emph{fuerza} a que
        vivan en el mismo espacio.
\end{enumerate}

\noindent\textbf{Analogía}: es como si dos personas midieran en
metros y otra persona midiera en pies. Si no comparten unidad,
no se pueden comparar. El DSC compartido es la "conversión a la
misma unidad".

\section{¿Qué hace \texttt{squeeze} y \texttt{contiguous}?}

\subsection{\texttt{squeeze(dim)}}

Elimina dimensiones de tamaño 1. Por ejemplo:
\begin{itemize}
  \item $(B, 1024, 1, 1) \xrightarrow{\text{squeeze}(-1)} (B,
        1024, 1)$.
  \item $(B, 1024, 1) \xrightarrow{\text{squeeze}(-1)} (B,
        1024)$.
\end{itemize}

\noindent El argumento $-1$ indica "la última dimensión". Útil
para preparar tensores para operaciones que esperan un número
específico de dims.

\subsection{\texttt{contiguous()}}

Algunos operaciones de PyTorch (especialmente \texttt{view},
\texttt{reshape}, slicing) requieren que el tensor sea
\emph{contiguo} en memoria. Si no lo es, se hace una copia
interna. Llamar \texttt{.contiguous()} fuerza la contigüidad
explícitamente.

\section{¿Por qué dimensión 1024?}

La dimensión del espacio conjunto ($1024$) es un compromiso:

\begin{itemize}
  \item \textbf{Suficientemente grande} para capturar la
        semántica de ambas modalidades.
  \item \textbf{Suficientemente pequeño} para:
        \begin{itemize}
          \item Similitud coseno estable numéricamente.
          \item Inferencia rápida (4 KB por embedding en FP32).
          \item Cabe en GPUs modestas.
        \end{itemize}
  \item \textbf{Coincide} con las unidades ocultas de la BiGRU,
        lo que permite compartir el último DSC sin mismatch de
        dimensiones.
\end{itemize}

\section{Código: Constructor \texttt{\_\_init\_\_}}

\begin{minted}[language=Python, numbers, firstnumber=13, linerange=13-44]{python}
class TextReIDNet(nn.Module):
    """La clase principal del modelo."""

    def __init__(self, configs: dict = None):
        # Llama al constructor de nn.Module (registra parámetros,
        # buffers, submódulos, etc.)
        super(TextReIDNet, self).__init__()

        # Validación: configs es obligatorio
        if configs is None:
            raise ValueError("`configs` parameter can not be None")

        # Guardar referencia a la configuración completa
        self.configs = configs

        # --- Rama visual ---
        # VisualNetwork = EfficientNet-B0 + wrapper
        self.visual_network = VisualNetwork()

        # --- Rama de lenguaje ---
        # GRULanguageNetwork = Embedding + Dropout + BiGRU
        self.language_network = GRULanguageNetwork(configs)

        # --- Capas de proyección al espacio conjunto ---
        # 1ra DSC: reduce canales de 1280 (EffNet) a 1024 (joint dim)
        self.visual_features_downscale = DepthwiseSeparableConv(1280, 1024)

        # Adaptive max pooling: colapsa (H, W) a (1, 1)
        self.adaptive_max_pooling = nn.AdaptiveMaxPool2d((1, 1))

        # 2da DSC: mantiene 1024 canales, proyecta al espacio conjunto
        # COMPARTIDA con la rama textual (ver text_embedding)
        self.depthwise_seperable_convolution = DepthwiseSeparableConv(
            1024, self.configs.feature_length
        )
\end{minted}

\section{Código: \texttt{forward}}

\begin{minted}[language=Python, numbers, firstnumber=45, linerange=45-59]{python}
    def forward(self, image, text_ids=None, text_length=None):
        """image: (B,3,H,W), text_ids: (B,100), text_length: (B,)
           returns: ((B,1024,1), (B,1024,1))"""
        # Calcular embedding visual (rama izquierda)
        visaual_embeddings = self.image_embedding(image)

        # Calcular embedding textual (rama derecha)
        textual_embeddings = self.text_embedding(text_ids, text_length)

        # Devolver ambos. train.py los aplana a (B, 1024) con view(-1).
        # Nota: hay un typo en el código: "visaual" en lugar de "visual"
        return visaual_embeddings, textual_embeddings
\end{minted}

\section{Código: \texttt{image\_embedding}}

\begin{minted}[language=Python, numbers, firstnumber=62, linerange=62-72]{python}
    def image_embedding(self, image=None):
        """image: (B, 3, H, W) batch de imágenes.
           returns: (B, 1024, 1) embeddings visuales."""

        # PASO 1: EfficientNet-B0
        # Entrada: (B, 3, 512, 512)
        # Salida: (B, 1280, 16, 16) <- última etapa
        visual_features = self.visual_network(image)

        # PASO 2: DSC 1280 -> 1024
        # Entrada: (B, 1280, 16, 16)
        # Salida:  (B, 1024, 16, 16)
        visual_features = self.visual_features_downscale(visual_features)

        # PASO 3: Adaptive Max Pooling
        # Entrada: (B, 1024, 16, 16)
        # Salida:  (B, 1024, 1, 1)
        visual_features = self.adaptive_max_pooling(visual_features)

        # PASO 4: DSC final 1024 -> 1024 (compartido con texto)
        # Entrada: (B, 1024, 1, 1)
        # Salida:  (B, 1024, 1, 1)
        visual_features = self.depthwise_seperable_convolution(visual_features)

        # PASO 5: squeeze + contiguous
        # .squeeze(-1) elimina la última dim (1) -> (B, 1024, 1)
        # .contiguous() asegura memoria contigua
        visual_features = visual_features.squeeze(-1).contiguous()
        return visual_features
\end{minted}

\section{Código: \texttt{text\_embedding}}

\begin{minted}[language=Python, numbers, firstnumber=75, linerange=75-89]{python}
    def text_embedding(self, text_ids=None, text_length=None):
        """text_ids: (B, 100), text_length: (B,).
           returns: (B, 1024, 1) embeddings textuales."""

        # PASO 1: Pasar por la red de lenguaje
        # Salida de GRULanguageNetwork:
        # Entrada:  (B, 100)         tokens BERT
        #   Embedding: (B, 100, 512)
        #   BiGRU:     (B, 100, 2048) bidireccional
        #   avg f+b:   (B, 100, 1024) promedio de fwd + bwd
        #   permute:   (B, 1024, 100, 1)
        textual_features = self.language_network(text_ids, text_length)

        # PASO 2: Max-pooling sobre la dim de tokens (dim=2)
        # Para cada canal (1024), tomar el valor máximo sobre los 100 tokens
        # Entrada: (B, 1024, 100, 1)
        # Salida:  (B, 1024, 1, 1)
        textual_features, _ = torch.max(textual_features, dim=2, keepdim=True)

        # PASO 3: DSC compartido con la rama visual
        textual_features = self.depthwise_seperable_convolution(textual_features)

        # PASO 4: squeeze + contiguous
        textual_features = textual_features.squeeze(-1).contiguous()
        return textual_features
\end{minted}

\section{Diagrama del flujo completo}

\begin{figure}[H]
\centering
\begin{tikzpicture}[
    font=\sffamily\footnotesize,
    node distance=0.7cm,
    box/.style={
        rectangle, draw=black, thick, rounded corners=2pt,
        minimum height=0.8cm, minimum width=2.2cm,
        align=center, fill=blue!10
    },
    arr/.style={-{Stealth[length=2.5mm,width=2mm]}, thick}
]
\node[box] (i) {Imagen\\$(B, 3, 512, 512)$};
\node[box, below=of i] (eff) {EffNet-B0\\$(B, 1280, 16, 16)$};
\node[box, below=of eff] (d1) {DSC\\$1280 \to 1024$};
\node[box, below=of d1] (p) {Pool $(1,1)$};
\node[box, below=of p] (d2) {DSC\\$1024 \to 1024$};
\node[box, below=of d2, fill=green!15] (ve) {$\mathbf{I}' (B, 1024, 1)$};

\node[box, right=3cm of i] (t) {Tokens\\$(B, 100)$};
\node[box, right=3cm of eff] (e) {Embedding\\$(B, 100, 512)$};
\node[box, right=3cm of d1] (g) {BiGRU\\$(B, 100, 2048)$};
\node[box, right=3cm of d2] (m) {Max$+$Perm\\$(B, 1024, 1, 1)$};

\draw[arr] (i) -- (eff);
\draw[arr] (eff) -- (d1);
\draw[arr] (d1) -- (p);
\draw[arr] (p) -- (d2);
\draw[arr] (d2) -- (ve);

\draw[arr] (t) -- (e);
\draw[arr] (e) -- (g);
\draw[arr] (g) -- (m);
\draw[arr] (m.east) -| ([yshift=0.4cm]ve.east) -- (ve.east);

\draw[arr, dashed] (ve) -- node[right] {cosine} +(2, 0);
\end{tikzpicture}
\caption{Flujo completo de tensores. Notar que el último DSC es
compartido entre ambas ramas.}
\label{fig:textreidnet-flow}
\end{figure}

\section{Inicialización de pesos}

\texttt{TextReIDNet} no llama ninguna función de inicialización.
Los pesos heredan:

\begin{itemize}
  \item \textbf{EfficientNet-B0}: preentrenado en ImageNet.
  \item \textbf{Embedding}: $\mathcal{N}(0, 1)$ por defecto de
        PyTorch.
  \item \textbf{BiGRU}: inicialización por defecto de PyTorch.
  \item \textbf{DSC}: Kaiming (\texttt{weights\_init\_kaiming}).
\end{itemize}

\section{Inferencia: extraer embeddings}

\begin{minted}[language=Python, basicstyle=\ttfamily\small, frame=single]{python}
import torch
from config import sys_configuration
from model.textreidnet import TextReIDNet

cfg = sys_configuration(dataset_source='huggingface')
model = TextReIDNet(cfg).to(cfg.device).eval()

# Modo evaluación (desactiva dropout)
with torch.no_grad():
    # Embedding de imagen
    img = torch.randn(2, 3, 512, 512).to(cfg.device)
    img_emb = model.image_embedding(img)        # (2, 1024, 1)
    img_emb = img_emb.view(img_emb.size(0), -1)  # (2, 1024)

    # Embedding de texto
    txt = torch.zeros(2, 100, dtype=torch.long).to(cfg.device)
    lng = torch.tensor([50, 30]).to(cfg.device)
    txt_emb = model.text_embedding(txt, lng)     # (2, 1024, 1)
    txt_emb = txt_emb.view(txt_emb.size(0), -1)  # (2, 1024)

    # Similitud coseno entre las 2 imágenes y los 2 textos
    img_emb_n = img_emb / img_emb.norm(dim=1, keepdim=True)
    txt_emb_n = txt_emb / txt_emb.norm(dim=1, keepdim=True)
    sim = torch.mm(img_emb_n, txt_emb_n.t())     # (2, 2)
    print(sim)
\end{minted}
